\documentclass[12pt]{article}

\usepackage[margin=1in]{geometry}
\usepackage{amsmath}
\usepackage{listings}
\usepackage{hyperref}
\usepackage{csquotes}
\usepackage[style=apa,backend=biber]{biblatex}
\DeclareLanguageMapping{american}{american-apa}
\addbibresource{references.bib}

\title{Working Title for the Mini-Research Paper}
\author{Author Name(s)}
\date{\today}

\begin{document}

\maketitle

\section*{Title}
State the final research title clearly and concisely.

\section*{Abstract}
Provide a concise summary of the research motivation, methods, findings, and conclusions.

\section{Introduction}
Introduce the problem context, motivation, and key research objectives.

\section{Preliminaries}
Define core concepts, notation, assumptions, and any prerequisite background.

\section{Methods and Applications}

\subsection{Model or Method Description}
Describe the mathematical model or algorithm. Include key equations such as:
\[
\frac{dN}{dt} = rN\left(1 - \frac{N}{K}\right)
\]

\subsection{Implementation and Code}
Explain the implementation choices, software structure, and reproducibility details.

\subsection{Results and Visualization}
Summarize generated tables/plots and describe how results are presented.

\subsection{Analysis and Discussion}
Interpret findings, discuss limitations, and compare outcomes with expectations.

\section{Conclusion}
Summarize the main contributions and list possible next steps.

\section*{References}
\printbibliography[heading=none]

\end{document}
